% TOP_PROBLEM: {"id": 30007445, "title": "FPRAS for cell-bounded contingency tables", "category_id": 15, "category": {"id": 15, "name": "computer_science", "display_name": "Computer Science", "description": "Computational complexity, algorithms, and theoretical CS.", "slug": "computer-science", "order_index": 15, "created_at": "2026-07-31T15:26:25.671Z"}, "status": "open"}
% FIRST_POSTED: 2026-10-06
% ATTEMPT_STATUS: solved
% OPENAI_CLAIM: {"schema_version":1,"source_repo":"https://github.com/openai/math","source_commit":"adc7f1241b42e322a6451854ab7e4b4c146bf78a","release_date":"2026-10-06","match":"full","resolution":"proved","independently_reviewed":false,"second_pass":"agreed","families":[{"family":"115","title":"Sampling and counting contingency tables with arbitrary margins","subject":"Theoretical computer science","manuscripts":[{"title":"An FPRAS for Cell-Bounded Contingency Tables","pdf_path":"preprints/An-FPRAS-for-Cell-Bounded-Contingency-Tables-September-24-2026/main.pdf","readme_path":"preprints/An-FPRAS-for-Cell-Bounded-Contingency-Tables-September-24-2026/README.md","date":"2026-09-24","theorem_ref":"preprints/An-FPRAS-for-Cell-Bounded-Contingency-Tables-September-24-2026/build/sections/introduction.tex#thm:main"}],"lean":{"doc_path":"lean/docs/115.md","comparators":["lean/ComparatorChallenges/ContingencyTables.json","lean/ComparatorChallenges/ContingencyTables.lean"],"declarations":[],"covers_main_theorem":true},"reasoning_trace":null}]}
\documentclass[11pt]{article}
\usepackage[margin=25mm]{geometry}
\usepackage{fontspec}
\setmainfont{Latin Modern Roman}
\usepackage{amsmath,amssymb,mathtools}
\usepackage{unicode-math}
\setmathfont{Latin Modern Math}
\usepackage{enumitem}
\usepackage{xurl,hyperref,bookmark}
\hypersetup{colorlinks=true,urlcolor=blue}
\setcounter{secnumdepth}{0}
\setlength{\parindent}{0pt}
\setlength{\parskip}{5pt}
\setlength{\emergencystretch}{3em}
\begin{document}
\section*{FPRAS for cell-bounded contingency tables}

\subsection{OpenAI's claim}
OpenAI presents the following claims relevant to this question, with manuscript numbers following paper order in Sources. In manuscript 1 (family 115; 2026-09-24), the asserted result is that binary-encoded contingency tables with individual cell bounds admit approximate counting with relative-error confidence guarantees and polynomial work on every random execution.

\subsection{Outline of the argument}
For manuscript 1 (family 115; 2026-09-24), softened integral signatures and balanced summation extend weights to defect states. A transport inequality controls the annealing chain, and bin-based rational weight estimation with bounded fair-bit simulation provides a finite implementation.

\subsection{Scope relative to this problem}
The result supplies the conclusion in this problem with its displayed hypotheses and quantifiers. The broader statements and applications in the manuscript are kept separate from the particular question above. The supporting result is manuscript 1 (family 115; 2026-09-24).

\subsection{Formal verification}
For manuscript 1 (family 115; 2026-09-24), OpenAI lists Lean statements for this assertion: binary-encoded contingency tables with individual cell bounds admit approximate counting with relative-error confidence guarantees and polynomial work on every random execution.

The selected statement includes the clause used for this problem.

The linked family Comparator checks are \texttt{ContingencyTables.lean}.

Their named declarations include \texttt{OAI.ContingencyTables.boundedSampling}, \texttt{OAI.ContingencyTables.counting}, \texttt{OAI.ContingencyTables.exactSampling}.

This site has not executed the Lean code. Described statement coverage is not an independent proof check.

\subsection{Status caveat}
The manuscripts were produced by an internal OpenAI model. OpenAI cautions that results without formalization can contain mistakes. This site has not independently checked the proofs; the label records the scope of OpenAI's claim.

\subsection{Sources}
\begin{itemize}
\item OpenAI, \emph{An FPRAS for Cell-Bounded Contingency Tables}, 2026-09-24; citation key \texttt{OAI:An-FPRAS-for-Cell-Bounded-Contingency-Tables-September-24-2026}. \href{https://github.com/openai/math/blob/adc7f1241b42e322a6451854ab7e4b4c146bf78a/preprints/An-FPRAS-for-Cell-Bounded-Contingency-Tables-September-24-2026/main.pdf}{PDF (GitHub)}; \href{https://github.com/openai/math/raw/adc7f1241b42e322a6451854ab7e4b4c146bf78a/preprints/An-FPRAS-for-Cell-Bounded-Contingency-Tables-September-24-2026/main.pdf}{PDF (direct)}; \href{https://github.com/openai/math/blob/adc7f1241b42e322a6451854ab7e4b4c146bf78a/preprints/An-FPRAS-for-Cell-Bounded-Contingency-Tables-September-24-2026/README.md}{Manuscript page and citation}; \href{https://github.com/openai/math/blob/main/preprints/An-FPRAS-for-Cell-Bounded-Contingency-Tables-September-24-2026/main.pdf}{Latest version}.
\item Family 115: \href{https://github.com/openai/math/blob/adc7f1241b42e322a6451854ab7e4b4c146bf78a/lean/docs/115.md}{Lean scope}.
\item \href{https://github.com/openai/math/blob/adc7f1241b42e322a6451854ab7e4b4c146bf78a/lean/ComparatorChallenges/ContingencyTables.json}{Comparator ContingencyTables.json}.
\item \href{https://github.com/openai/math/blob/adc7f1241b42e322a6451854ab7e4b4c146bf78a/lean/ComparatorChallenges/ContingencyTables.lean}{Comparator ContingencyTables.lean}.
\item \href{https://github.com/openai/math/blob/adc7f1241b42e322a6451854ab7e4b4c146bf78a/CONTENTS.md}{OpenAI catalogue entry}.
\end{itemize}
\end{document}
